\chapter{결론 및 향후 연구 방향}
\section{결론}
본 과제는 고전 컴퓨터 비전과 역기구학을 이용한 SO-101 블록 이동 제어를 구현하고 실기로 부분 수행하였다. 적재는 PLACE 전략과 접촉 감지의 구현 단계이며, 타워 자세 등록과 통합 실기 검증이 남아 있다. 인식, 파지 후보 계산, 위치/부하 확인, 재시도와 슬롯 배치를 연결하여 동작 결과에 따라 다음 상태를 결정하는 구조를 구성하였다.

외부 180초 제한의 전체 미션 한 회에서는 파지 후보 14회의 실행 시작, HELD 4회, EMPTY 9회와 판정 전 중단 1회를 기록하였다. 네 번의 슬롯 해제 후 마지막 블록이 남아 미션은 미완료였으며, 해제 수를 인정 블록 수로 환산하지 않았다. 별도 단일 블록 시험에서는 첫 위치의 파지와 운반에 성공했지만 안쪽 위치의 기본 시도와 보정 재시도는 빈손으로 판정되었다. 최종 CV+IK의 동일 조건 반복 성공률과 ACT 대비 성능 차이는 측정하지 못하였다.

실행 로그는 노란 블록의 후보 재시도와 빨간 블록의 재선택 후 파지 판정 회복을 보여준다. 반면 빨간 블록의 첫 PICK에 51.5초가 소요되어 재시도의 시간 비용도 확인하였다. 좌표 정합은 RMS 8.18\,mm, 최대 LOO 26.64\,mm로 개발 목표에 미달하였다. 본 연구의 주요 기여는 기존 라이브러리의 영상 처리와 기구학 기능을 과제에 맞는 선택, 계획, 검증 및 복구 흐름으로 통합하고, 실패 단계를 실행 기록으로 분석한 데 있다. 요구된 미션의 안정적인 완수를 입증하려면 다음의 검증과 개선이 필요하다.

\section{향후 연구 방향}
\subsection{동일 조건 반복 평가와 시간 예산 개선}
우선 한 초기 배치를 사진으로 고정하고 코드, 실행 설정과 캘리브레이션을 유지한 상태에서 전체 미션을 3회 반복할 계획이다. 각 시행에 대해 사람이 판정한 인정 블록 수, 시스템 완료 상태, 제한시간 종료와 실제 정리 동작 시간을 따로 기록한다. 파지 영상과 센서 판정을 연결하고 첫 후보, PICK 내부 재시도 및 대상 재선택을 구분한다. 이 소규모 시험은 고정 배치의 반복성 평가이며, 무작위 배치의 일반화 평가는 별도로 수행한다.

현재의 실패 기록을 바탕으로 후보별 시간 상한과 다음 블록으로의 전환 시점을 우선 검토한다. 성능 측정 중에는 설정을 변경하지 않고, 이후 같은 배치에서 재시도 정책을 변경한 결과를 비교한다. 실제 턱과 블록의 상대 위치 측정, 파지 직전 영상 보정, 검출 누락과 물리적인 완료 조건의 차이도 후속 개선 대상으로 둔다.

\subsection{적재 높이별 동작 개선과 실기 평가}
\label{sec:stack-extension}
적재는 타워 접근과 하강 자세를 등록한 뒤 높이별 도달성, 관절 한계와 그리퍼 개방 공간부터 확인한다. 접촉 감지와 되돌림 후 해제가 아래 블록을 밀거나 타워를 흔드는지 2단과 3단에서 평가하고, 이후 다섯 블록과 5초 유지 조건으로 확대한다. 각 단계의 시도 횟수, 실패 단계와 유지시간을 기록하며 현재는 낮은 단수의 실기 성공도 확인되지 않았다.

평면 변환은 블록 한 개의 윗면을 기준으로 하므로 타워가 높아졌을 때 같은 변환을 그대로 적용할 수 없다. 초기에는 고정한 타워 위치와 기록 자세를 사용하고, 높이 변화에 따른 관측 오차와 적재물 재선택 문제를 별도로 평가한다.

\subsection{규칙 기반 실행의 학습 데이터 활용}
\label{sec:dataset-extension}
미션의 반복 수행과 결과 판정이 안정화되면 CV+IK 실행 중 영상, 관절 상태와 전달 명령을 같은 시간 기준으로 기록하는 확장을 검토한다. 초기에는 사람이 블록을 재배치하고, 단일 블록 선택부터 배치 후 재관측까지를 한 에피소드로 구성한다. 내부 재시도, 실패, 시간 초과와 사람의 확인 결과를 함께 저장하며, 현재의 사진 및 텍스트 기록에 동기화와 에피소드 구성 기능을 추가해야 한다.

저장 형식은 사용할 LeRobot 버전과 호환되도록 선택한다\cite{lerobot-dataset}. 모방학습이나 SmolVLA에 활용하려면 관측과 행동 표현, 작업 지시 및 센서 라벨의 신뢰도를 확인해야 한다\cite{smolvla}. 강화학습용으로 확장할 때에는 다음 관측, 보상과 종료 조건도 정의하고, 파지 라벨을 전체 배치 성공의 보상으로 그대로 대체하지 않는다. 유효 에피소드 수, 기록 누락과 사람의 작업시간으로 수집 방식을 평가하며, 데이터 수집 자동화와 학습 성능 개선은 현재 실증한 성과에 포함하지 않는다.
\FloatBarrier
